Large language models whose weights exceed the memory of a graphics processor can be served by keeping some layers resident and streaming the others from host-attached memory, such as page-locked host memory or a Compute Express Link expander, through a small ring of on-device buffers. Common tools fill device memory in model order. We show that the position of the resident layers matters and answer the placement question exactly for layers of equal cost. We model one inference iteration as a max-plus linear system and prove that its iteration time is the maximum mean of a family of tilings of the layer sequence. Three results follow. Resident layers are useless when a layer computes at least as long as it takes to copy. Otherwise, spreading the resident layers evenly attains the lower bound set by total computation and total copying with three buffers, and with two when a copy takes at least twice as long as a layer's computation. A contiguous block of resident layers, as produced by model-order filling, has an exact closed-form iteration time that approaches twice the optimum. On a Tesla T4 with page-locked copies at \nHwBw~GB/s and real Llama-3.2-3B and Llama-3.1-8B layer kernels, the model, fed with per-layer times measured on the same device, reproduces \nAccN{} measured iteration times with a mean error of \nAccMape\% and orders \nAccRankExact\% of distinguishable placement pairs correctly. Even spreading is up to \nMeasPenaltyMax{} times faster than a contiguous block.
